\section{Summary}
Scientific-ML researcher with a physics-first background in equivariant graph networks, automatic differentiation, first-principles modeling, and GPU/HPC implementation.
\section{Technical Skills}
\SkillLine{ML}{JAX, PyTorch, PyTorch Geometric, e3nn, equivariant GNNs, automatic differentiation, NumPy, SciPy}
\SkillLine{Scientific compute}{Python, C++, MPI, OpenMP, CUDA (cuBLAS, cuTENSOR), ROCm (rocBLAS, rocFFT), PETSc, SLEPc, DOLFINx}
\SkillLine{Platforms}{Frontier, Perlmutter, Frontera, Docker, Kubernetes, Linux}
\section{Experience}
\Role{PhD Researcher}{Yale University}{Electrical Engineering}{2021--present}
\begin{itemize}
\item Build DeePH-style computational graphs for LiF and double-perovskite systems; train equivariant models in PyTorch Geometric and use automatic differentiation to obtain force and optical-property quantities.
\item Connect ML models to first-principles exciton and many-body theory, ensuring datasets, losses, and evaluations map to physical observables.
\item Build and debug distributed GPU workflows in Python/C++ using MPI, OpenMP, CUDA/ROCm, PETSc, and SLEPc on national HPC systems.
\item Prototype retrieval-driven scientific assistants that extract structured simulation inputs from papers and technical documentation.
\end{itemize}
\Role{Software and Embedded Systems Engineer}{Probe Technologies / Weatherford Wireline Services}{Houston, TX}{2019--2021}
\begin{itemize}
\item Built CUDA-enabled waveform-processing and visualization tools, adding production experience with data pipelines and performance constraints.
\item Implemented SQL-backed ASP.NET services and C\# tools for remote engineering-test workflows.
\end{itemize}
\section{Selected Projects}
\Project{EGNN QM9 force model}{Trained an equivariant GNN and differentiated predicted molecular energies to obtain forces.}
\Project{Materials optical-property models}{Applied graph networks and automatic differentiation to accelerate costly calculation stages; roughly 100x speedup in parts of the workflow.}
\Project{Scientific retrieval pipeline}{Structured software documentation and paper content into starter artifacts for computational-science workflows.}
\section{Education}
\Role{PhD, Electrical Engineering}{Yale University}{Computational materials and scientific ML}{2021--present}
\Role{MPhil and MS, Electrical Engineering}{Yale University}{}{2023}
\Role{MEng and BS, Electrical and Computer Engineering}{Cornell University}{}{2013--2018}
\ifcv
\section{Research Record}
Co-author, \textit{Journal of the American Chemical Society}, 2025; APS presentations, 2024--2026.\\
\fi
